\documentclass[10pt,letterpaper]{article}
\usepackage{cornell-notes}

\title{Chapter 5: The Network Layer}
\author{}
\date{}

\setDocTitle{Chapter 5: The Network Layer}
\setDocSubtitle{Computer Networks}
\setDocVersion{v1.0}
\setDocStatus{Study Companion}
\setDocOwner{Computer Networks Cornell Notes}
\setDocAuthor{}
\setDocDate{\today}
\setDocSummary{Cornell-format study and review notes for Chapter 5: The Network Layer.}

\setCornellCollection{Computer Networks Cornell Notes}
\setCornellUnitType{Chapter}
\setCornellUnitNumber{5}
\setCornellUnitTitle{The Network Layer}
\setCornellCompanionLabel{Study and review companion}

\hypersetup{pdftitle={Chapter 5: The Network Layer},pdfsubject={Cornell notes},pdfauthor={}}

\begin{document}
\maketitle
\notesection{Chapter overview}
\begin{overviewbox}
\textbf{Learning objectives}\begin{itemize}
  \item Distinguish datagram forwarding from virtual-circuit forwarding and routing from forwarding.
  \item Trace shortest-path, flooding, distance-vector, link-state, hierarchical, group, and mobile routing approaches.
  \item Explain congestion control and quality-of-service mechanisms from admission through scheduling.
  \item Describe tunneling, fragmentation, and routing across heterogeneous networks.
  \item Interpret IPv4, IPv6, control protocols, MPLS, OSPF, and BGP roles.
\end{itemize}
\textbf{Prerequisite concepts}\quad Graphs and shortest paths, packet switching, link-layer framing, probability, and queueing intuition.\par
\textbf{Key themes}\quad Forwarding state; route computation; scale; congestion; service differentiation; heterogeneous interconnection; address architecture.\par
\textbf{Named mechanisms and standards}\quad Dijkstra, flooding, distance vector, link state, multicast trees, traffic shaping, fair queueing, IntServ, DiffServ, IPv4, CIDR, NAT, IPv6, ICMP, ARP, DHCP, MPLS, OSPF, BGP, multicast, and Mobile IP.
\end{overviewbox}

\notesection{Cornell cue-and-note record}

\begin{longtable}{@{}L{0.245\textwidth}|L{0.695\textwidth}@{}}
\rowcolor{CNavy}\color{white}\bfseries Cue / recall prompt &
\color{white}\bfseries Notes \\
\endfirsthead
\rowcolor{CNavy}\color{white}\bfseries Cue / recall prompt &
\color{white}\bfseries Notes (continued) \\
\endhead
\hline
\endfoot
\cellcolor{CCue}\textbf{5.1 Design issues} & \begin{minipage}[t]{\linewidth}\raggedright Forwarding handles an individual packet using current state; routing computes or distributes the state used for forwarding.\par\begin{itemize}\item \textbf{5.1.1 Store-and-forward switching} A router receives a packet, checks it, selects an outgoing interface, and queues it for transmission. Each hop contributes processing, queueing, transmission, and propagation delay.\item \textbf{5.1.2 Service to transport} A network layer may expose connectionless datagram service or connection-oriented virtual circuits, with differing guarantees and state placement.\item \textbf{5.1.3 Connectionless implementation} Each packet carries a destination and is forwarded independently from a table. Route changes can cause packets in one flow to take different paths or arrive out of order.\item \textbf{5.1.4 Connection-oriented implementation} Setup installs virtual-circuit identifiers and forwarding entries along a path; packets carry a short local identifier until teardown.\item \textbf{5.1.5 Virtual circuits versus datagrams} Virtual circuits move complexity and failure sensitivity into connection state; datagrams carry fuller addresses and let routers remain less dependent on per-flow setup.\end{itemize}\end{minipage} \\ \hline
\cellcolor{CCue}\textbf{5.2 Routing algorithms} & \begin{minipage}[t]{\linewidth}\raggedright \begin{itemize}\item \textbf{5.2.1 Optimality principle} If router $J$ lies on an optimal path from $I$ to $K$, the optimal path from $J$ to $K$ is the remainder of that path. Destination-rooted optimal routes form a sink tree when costs are stable.\item \textbf{5.2.2 Shortest path} Dijkstra's algorithm repeatedly makes the tentative node with least path cost permanent and relaxes its outgoing edges. Correctness assumes suitable nonnegative link costs.\item \textbf{5.2.3 Flooding} A router sends a packet on every outgoing link except the arrival link. Hop limits, sequence numbers, or selective direction control duplication; flooding is robust but expensive.\item \textbf{5.2.4 Distance vector} Routers exchange destination cost vectors with neighbors and apply the Bellman-Ford recurrence. Bad news can propagate slowly, producing routing loops and count-to-infinity behavior.\item \textbf{5.2.5 Link state} Each router discovers neighbors and costs, floods a link-state packet, builds a topology map, and computes shortest paths locally. Sequence numbers and aging control stale advertisements.\item \textbf{5.2.6 Hierarchical routing} Grouping routers into regions reduces table and update size, at the cost of routes that may be less than globally optimal.\item \textbf{5.2.7 Broadcast routing} Methods include individual unicasts, flooding, multidestination packets, sink-tree forwarding, and reverse-path forwarding.\item \textbf{5.2.8 Multicast routing} Delivery to group members prunes or constructs trees so packets reach subscribed receivers without traversing every link.\item \textbf{5.2.9 Anycast routing} One group address represents multiple instances; ordinary routing delivers a packet to a nearby or lowest-cost member rather than all members.\item \textbf{5.2.10 Mobile-host routing} A stable home identity is reconciled with a changing attachment point through registration and forwarding or tunneling.\item \textbf{5.2.11 Ad hoc routing} With no fixed infrastructure and changing radio neighbors, routes must be discovered and maintained dynamically; proactive and on-demand strategies trade overhead for discovery delay.\end{itemize}\end{minipage} \\ \hline
\cellcolor{CCue}\textbf{5.3 Congestion control} & \begin{minipage}[t]{\linewidth}\raggedright Congestion is sustained overload within the network; it differs from a receiver's inability to consume data.\par\begin{itemize}\item \textbf{5.3.1 Approaches} Open-loop prevention designs policies in advance; closed-loop control observes congestion, feeds information back, and adjusts offered load. Control may act hop by hop or end to end.\item \textbf{5.3.2 Traffic-aware routing} Routing can shift traffic away from overloaded links, but rapid reactions may move congestion or oscillate unless load and update dynamics are controlled.\item \textbf{5.3.3 Admission control} A new flow is accepted only if the network expects to meet existing and requested resource commitments.\item \textbf{5.3.4 Traffic throttling} Routers can signal congestion explicitly or indirectly; sources reduce load in response. Hop-by-hop backpressure reacts locally but propagates upstream.\item \textbf{5.3.5 Load shedding} When buffers and control are insufficient, routers discard packets. The discard policy can consider application value and protocol response rather than dropping blindly.\end{itemize}\end{minipage} \\ \hline
\cellcolor{CCue}\textbf{5.4 Quality of service} & \begin{minipage}[t]{\linewidth}\raggedright \begin{itemize}\item \textbf{5.4.1 Application requirements} Applications differ in bandwidth, delay, jitter, and loss tolerance. Elastic file transfer and real-time media therefore value different service properties.\item \textbf{5.4.2 Traffic shaping} Leaky-bucket shaping emits at a controlled rate; token-bucket shaping permits bursts while bounding long-term volume. Shaping changes traffic presented to the network.\item \textbf{5.4.3 Packet scheduling} FIFO is simple; priority can starve low classes; fair and weighted fair queueing approximate controlled sharing among flows or classes.\item \textbf{5.4.4 Admission control} A flow specification is checked against route resources before a guarantee is accepted; policing then verifies conformance.\item \textbf{5.4.5 Integrated services} Per-flow reservation and signaling can support quantitative guarantees, but per-flow state and setup burden limit scale.\item \textbf{5.4.6 Differentiated services} Packets are marked into behavior classes, and routers apply configured per-hop treatment. Aggregation improves scalability but offers less individualized control.\end{itemize}\end{minipage} \\ \hline
\cellcolor{CCue}\textbf{5.5 Internetworking} & \begin{minipage}[t]{\linewidth}\raggedright \begin{itemize}\item \textbf{5.5.1 Network differences} Networks may differ in addressing, maximum packet size, service model, routing, quality, security, and timing.\item \textbf{5.5.2 Connecting networks} Repeaters and bridges extend lower-layer technologies; routers connect at the network layer; translation above the network layer is needed when semantics differ fundamentally.\item \textbf{5.5.3 Tunneling} A packet for one protocol is encapsulated inside another protocol while crossing an intervening network, then decapsulated at the tunnel endpoint.\item \textbf{5.5.4 Internetwork routing} Routing across administrative and technological boundaries uses abstraction, hierarchy, and policy; the route selected need not be the shortest physical path.\item \textbf{5.5.5 Fragmentation} A packet too large for a link may be divided into fragments and later reassembled. Transparent and nontransparent schemes differ in where reassembly occurs; fragmentation adds overhead and loss sensitivity.\end{itemize}\end{minipage} \\ \hline
\cellcolor{CCue}\textbf{5.6 Internet network layer} & \begin{minipage}[t]{\linewidth}\raggedright \begin{itemize}\item \textbf{5.6.1 IPv4} The header carries version, length, service information, total length, fragmentation fields, lifetime, upper-layer protocol, checksum, and addresses. Routers decrement lifetime and recompute the header checksum.\item \textbf{5.6.2 IP addresses} Hierarchical prefixes support subnetting and classless interdomain routing. Longest-prefix match selects a forwarding entry. DHCP configures hosts; NAT maps private internal addresses and transport identifiers to external reachability.\item \textbf{5.6.3 IPv6} A larger address space, simplified base header, extension headers, flow labeling, and improved autoconfiguration address IPv4 limitations. Transition requires coexistence and tunneling strategies.\item \textbf{5.6.4 Internet control protocols} ICMP reports network conditions and errors; ARP maps IPv4 addresses to local-link addresses; DHCP assigns configuration. Each solves a control problem adjacent to IP forwarding.\item \textbf{5.6.5 Label switching and MPLS} Ingress classification assigns a short label; intermediate routers swap labels along a label-switched path. MPLS can steer traffic and separate forwarding classes below ordinary IP lookup.\item \textbf{5.6.6 OSPF} An interior link-state protocol floods topology information within areas, computes shortest paths, supports multiple metrics and equal-cost paths, and uses hierarchy around a backbone area.\item \textbf{5.6.7 BGP} An interdomain path-vector protocol advertises reachable prefixes with path and policy attributes. Autonomous systems choose routes according to policy as well as reachability.\item \textbf{5.6.8 Internet multicasting} Hosts signal group membership locally and multicast routing builds distribution state so one sent packet can reach multiple receivers efficiently.\item \textbf{5.6.9 Mobile IP} A mobile node retains a home address, obtains reachability at a visited network, registers it, and receives tunneled packets through a mobility agent or corresponding mechanism.\end{itemize}\end{minipage} \\ \hline
\cellcolor{CCue}\textbf{5.7 Summary} & \begin{minipage}[t]{\linewidth}\raggedright The network layer forwards packets across multiple links. Routing constructs scalable paths, congestion control protects shared capacity, QoS differentiates service, and IP plus its control and routing protocols interconnect heterogeneous networks.\end{minipage} \\ \hline
\end{longtable}

\begin{legacybox}Classful addressing, some IPv4 mobility-agent terminology, and the presented NAT and transition practices are retained as historical context. They illuminate deployed IPv4 constraints but are not a complete current architecture guide.\end{legacybox}

\notesection{Terminology and notation}

\begin{longtable}{@{}L{0.245\textwidth}|L{0.695\textwidth}@{}}
\rowcolor{CNavy}\color{white}\bfseries Cue / recall prompt &
\color{white}\bfseries Notes \\
\endfirsthead
\rowcolor{CNavy}\color{white}\bfseries Cue / recall prompt &
\color{white}\bfseries Notes (continued) \\
\endhead
\hline
\endfoot
\cellcolor{CCue}\textbf{Forwarding} & \begin{minipage}[t]{\linewidth}\raggedright Per-packet action that selects an output using a forwarding table.\end{minipage} \\ \hline
\cellcolor{CCue}\textbf{Routing} & \begin{minipage}[t]{\linewidth}\raggedright Process that computes, exchanges, or applies information used to build forwarding state.\end{minipage} \\ \hline
\cellcolor{CCue}\textbf{Datagram} & \begin{minipage}[t]{\linewidth}\raggedright Independently forwarded network-layer packet in a connectionless service.\end{minipage} \\ \hline
\cellcolor{CCue}\textbf{Virtual circuit} & \begin{minipage}[t]{\linewidth}\raggedright Pre-established path represented by per-hop identifiers and forwarding state.\end{minipage} \\ \hline
\cellcolor{CCue}\textbf{Congestion} & \begin{minipage}[t]{\linewidth}\raggedright Persistent demand beyond network resources, causing queue growth, delay, and loss.\end{minipage} \\ \hline
\cellcolor{CCue}\textbf{Jitter} & \begin{minipage}[t]{\linewidth}\raggedright Variation in packet delay.\end{minipage} \\ \hline
\cellcolor{CCue}\textbf{Prefix} & \begin{minipage}[t]{\linewidth}\raggedright Leading address bits identifying an aggregate network range.\end{minipage} \\ \hline
\cellcolor{CCue}\textbf{Longest-prefix match} & \begin{minipage}[t]{\linewidth}\raggedright Selection of the matching route with the greatest prefix length.\end{minipage} \\ \hline
\cellcolor{CCue}\textbf{Autonomous system} & \begin{minipage}[t]{\linewidth}\raggedright Administratively controlled routing domain represented in interdomain routing.\end{minipage} \\ \hline
\cellcolor{CCue}\textbf{Tunnel} & \begin{minipage}[t]{\linewidth}\raggedright Logical link created by encapsulating one packet format inside another.\end{minipage} \\ \hline
\end{longtable}

\notesection{Comparison matrix}

\begin{longtable}{@{}L{0.313\textwidth}L{0.313\textwidth}L{0.313\textwidth}@{}}
\toprule
\textbf{Approach} & \textbf{State / information} & \textbf{Characteristic tradeoff} \\ \midrule
Distance vector & Neighbor-reported destination costs & Simple exchange; slow bad-news convergence \\
Link state & Flooded topology and local shortest-path computation & Faster, richer view; more database and flood machinery \\
Integrated services & Per-flow reservation state & Fine-grained guarantees; scaling burden \\
Differentiated services & Aggregated class marking and per-hop behavior & Scalable classes; weaker per-flow specificity \\
Datagram network & Destination-based forwarding & Flexible failure response; larger per-packet addressing \\
Virtual-circuit network & Setup and per-connection forwarding entries & Short labels and control; path-state sensitivity \\
\bottomrule
\end{longtable}
\notesection{Chapter synthesis}

\begin{overviewbox}\textbf{Cause-and-effect chain.} Multiple links require forwarding; forwarding tables require routing; scale requires hierarchy and aggregation; shared queues create congestion; application diversity motivates scheduling and service differentiation; heterogeneous technologies require a common internetwork protocol.\par
\textbf{Common confusions.} Routing computes paths while forwarding handles packets; congestion control protects the network while flow control protects a receiver; ARP is local address resolution rather than a routing protocol; and BGP policy can override shortest-path intuition.\par
\textbf{Derived insight.} Network-layer scalability repeatedly exchanges fine-grained optimality for aggregation: hierarchical routes, prefixes, areas, autonomous systems, and service classes all compress state.\end{overviewbox}

\notesection{Retrieval practice}

\begin{enumerate}
  \item Trace how a datagram and a virtual-circuit packet are forwarded differently.
  \item Apply one iteration of Dijkstra relaxation to a small weighted graph.
  \item Why can distance-vector routing count to infinity after a failure?
  \item List the information a router must obtain before running a link-state computation.
  \item Compare broadcast, multicast, and anycast delivery goals.
  \item Distinguish congestion control from end-to-end flow control.
  \item How do leaky-bucket and token-bucket shaping differ for bursts?
  \item Why does weighted fair queueing provide more isolation than FIFO?
  \item Compare IntServ and DiffServ state and scaling behavior.
  \item Trace a tunneled packet through entry, transit, and exit networks.
  \item Why does fragmentation amplify the effect of losing one fragment?
  \item Perform a longest-prefix-match choice among overlapping prefixes.
  \item Explain how OSPF and BGP serve different administrative scopes.
  \item Why can a valid BGP route be non-shortest in physical distance?
\end{enumerate}
\begin{CornellSummaryBox}{Rapid-review Cornell summary}Forwarding moves packets; routing builds the state that forwarding uses. Datagram and virtual-circuit designs place state differently. Congestion mechanisms manage overload; QoS mechanisms shape, schedule, reserve, or classify traffic. IP provides common addressing and delivery, while OSPF, BGP, control protocols, and MPLS coordinate reachability and paths.\end{CornellSummaryBox}
\end{document}
